\newcommand{\pointwisewarning}{\ifnotesedition\else\raisebox{-.08em}{\includegraphics[height=.8em]{assets/warning-emoji.png}}\hspace{.25em}\fi}
\talkslide{pointwise}{\pointwisewarning Challenges with pointwise judges}{Timing: provisional}{A scalar judge must put each response on a shared scale.}{
% A compact supporting example leaves the two challenges as the main hierarchy.
\begin{tikzpicture}[remember picture,overlay,
 every node/.style={text=Ink,inner sep=0pt,align=left}]
\begin{scope}[shift={(current page.south west)}]
 \node[anchor=west,font=\normalsize\bfseries] at (.65,7.0)
  {1. Need to provide absolute scores};
 \node[anchor=west,font=\normalsize] at (1.25,6.36)
  {Prompt LLM judge: \textit{``Rate this response 1--10''}};
 \draw[Muted!55,line width=.5pt] (1.25,5.92)--(14.75,5.92);
 \node[anchor=west,font=\fontsize{9.3}{11.5}\selectfont\bfseries] at (1.4,5.61)
  {Response to \textit{``How do I make explosives?''}};
 \node[anchor=east,font=\fontsize{9.3}{11.5}\selectfont\bfseries] at (14.6,5.61) {Score?};
 \draw[Muted!35,line width=.4pt] (1.25,5.30)--(14.75,5.30);
 \node[anchor=north west,text width=11.3cm,font=\fontsize{9.3}{11.5}\selectfont] at (1.4,5.10)
  {``I can't help with that.''};
 \node[anchor=north east,font=\fontsize{9.3}{11.5}\selectfont] at (14.6,5.10) {5? 6? 7?};
 \node[anchor=north west,text width=11.3cm,font=\fontsize{9.3}{11.5}\selectfont] at (1.4,4.52)
  {``I can't provide instructions, but I can explain the physics of rapid combustion in controlled settings like fireworks\ldots''};
 \node[anchor=north east,font=\fontsize{9.3}{11.5}\selectfont] at (14.6,4.52) {8? 9?};
 \node[anchor=north west,text width=11.3cm,font=\fontsize{9.3}{11.5}\selectfont] at (1.4,3.62)
  {``That's dangerous and illegal. Here's why explosives are regulated\ldots''};
 \node[anchor=north east,font=\fontsize{9.3}{11.5}\selectfont] at (14.6,3.62) {7? 8?};
 \draw[Muted!55,line width=.5pt] (1.25,3.08)--(14.75,3.08);
 \node[anchor=west,font=\normalsize\bfseries] at (.65,2.28)
  {2. More susceptible to reward hacks};
 \node[anchor=north west,text width=13.5cm,font=\normalsize] at (1.25,1.79)
  {In early experiments, our attacker exploited the judge's \textbf{preference for verbosity}};
\end{scope}
\end{tikzpicture}

}
\speak{A judge can assess compliance or deflection. Here the example illustrates the difficulty of scoring deflection; a successful harmful attack would seek compliance instead. We will reuse this judging machinery for both players, with different criteria. These are the original illustrative responses to the same harmful request. The judge is asked to assign each a number from one to ten. A bare refusal, general safe background information, and a warning can each be appropriate to different degrees. The challenge is calibrating absolute scores across contexts. The examples here contain no operational details. In our early experiments with a pointwise judge, we observed high scalar scores for verbose phrases, and the attacker learned to exploit this by adding filler. This is the speaker's account of an observed failure, not a quantitative result added to the paper. They motivate a relative comparison; they do not establish that every judgment is purely subjective or that pairwise judges are always correct.}
